\chapter{77}



Sie knieten sich gegenüber, zwischen ihnen der eiserne Pumparm an der
Seite des Duschtanks. Louis schien der Moment geeignet, mehr über Joh
und Manuela zu erfahren. Die beiden hatten es ihm wirklich angetan. Fast
zu schön, wie die miteinander umgingen. Da musste doch irgendwo ein
Haken sein. Der Gedanke schlich sich immer wieder an.
Joh suchte sich die Werkzeuge aus der Kiste zusammen. Dann fasste sich
Louis ein Herz.

«Joh, sag mal, also, ich hätte da noch gern etwas von dir gewusst.»

Joh schaute hoch.

«Ja, nur zu.»

«Es geht um Manuela und dich. Du hast davon erzählt, wie ihr euch
kennen lerntet, dass ihr es langsam angingt, aber, ganz am Anfang, wie hast du sie von dir
überzeugt?»

Jetzt lag die Frage da. Joh schien sie zu gefallen. 

«Du meinst, wie ich eine Frau wie sie erobern konnte? Das meintest du
doch, oder?»

Johs Augen lachten. Er hatte Louis ertappt.

«Na, ja, auch.»

Bevor Joh an der Pumpe herumhantierte, legte er Louis die Hand auf den Arm. 
Er kniff die Augen zusammen und grinste.

«Okay, da müsste ich aber etwas ausholen. Magst du zuhören?»

Louis nickte nachdrücklich.

«Nun gut, also weisst du, als junger Mann hab ich mich mit der Liebe
schwer getan. Als Kind lief alles noch ganz gut, so galant, wie ich
Frauen und Mädchen gegenüber war, mon Dieu,» Joh verwarf die Hände, «na
ja, als Sohn eines Franzosen! Als Jugendlicher versuchte ich jegliche Ablehnung zu vermeiden - 
irgendeinmal platzte das Ganze. Ich habe dann lange selbstmitleidig in mir geschmort, bis ich wieder hochkam.»

Joh deutete auf den Boden neben Louis.

«Gibst du mir mal bitte die Zange?»

Louis reichte sie ihm. Bei einer Aussage war er hängen geblieben.

«Sohn eines Franzosen?»

«Ach, Franzosen lehrten damals die Söhne, Frauen gegenüber zuvorkommend
und galant zu sein. Mein Vater hatte es intus und ich imitierte ihn.
Damals wusste ich nicht, dass er ein Schürzenjäger war und wie sehr
meine Mutter darunter litt. Ich wollte ihm einfach nur gefallen. Schon
mit sechs, sieben Jahren verstand ich das Spiel - ah, aua, shit.»

Mit einem scheppernden Geräusch war der Arm der Pumpe nach unten
geschnellt und hatte Johs Daumen eingeklemmt. Louis hob blitzschnell die
Stange hoch. Und Joh zog seine Hand weg.

«Ah, Mensch, so ein Mist.»

Joh hielt hielt sich den Daumen. Louis berührte ihn an der Schulter.

«Soll ich den Verbandskasten holen?»

«Nein, nein, lass mal, geht schon wieder.»

Joh atmete durch. 

«Hm, wo waren wir? - ach ja, dann kam mein Untergang, ich verliebte mich so
richtig in ein Mädchen meiner Klasse und jetzt, wo es drauf ankam, griff
keine meiner Strategien,» er biss die Zähne zusammen und versuchte eine
verrostete Schraube zu lösen, «von mir als dem Süssen, Netten wollte sie
nichts wissen und bald schon war sie mit einem dieser frechen Typen
zusammen.»

Mit einem Ruck war es geschafft. Die Schraube löste sich. 

«Weisst du, mit meinen wahren Gefühlen habe ich mich schliesslich 
schwer verkracht - das war wohl reiner Selbstschutz, erkannte ich erst viel später.»

Joh schaute auf den Schraubenschlüssel in seiner Hand.
Louis konnte kaum glauben, wie offenherzig Joh erzählte. Es verwirrte ihn, dass er dabei komplett ruhig und zufrieden wirkte. Vielleicht sogar eine Art geheimnisvollen Glücks ausstrahlte?

«Es war ein echtes Ringen,» fuhr Joh fort, lachte nun, knapp überlaut, suchte in
der Werkzeugkiste nach einer passenden Schraube und redete weiter, «und
dann, ja, dann tauchte viel, viel später, nein noch später, Manuela in
der Leichenhalle auf. Davon habe ich dir ja erzählt.»

«Ja, genau. Aber, aber, wie lebtest du denn die zehn Jahre vor Manuela? Ich mein, du hast
Kinder, nicht?»

«Ja, ja, ich habe Zwillinge mit einer wirklich ausserordentlichen Frau.
Heute verstehe ich, warum das nicht klappen konnte.»

Louis sah das Familienbild an der Magnetwand vor sich. Daher also kam
die Ähnlichkeit der beiden jungen Leute, die mit Joh und Manuela darauf
posierten. Joh schob wieder die Brille zurecht und sprach weiter.

«Wir ersetzen die Schraube. Aber erstmal brauche ich das Schmierfett.
Schau, das da.»

Er deutete auf eine Tube. Louis schraubte den Deckel ab.

«Bitte.»

«Danke, ich musste fast fünfzig werden, um auf Manuela zu treffen. Nach
unserer ersten Begegnung in der Forensik verabredeten wir uns exzessiv
oft und redeten uns um Kopf und Kragen. Nächtelang. Ich weiss heute noch
nicht, wie es kam, dass ich Manuela \emph{so viel Johannes} auf den
Tisch legte.»

«Mutig,» fiel ihm Louis ins Wort. Er konnte sich nicht zurückhalten.

«Ich glaub nicht, Ciro, viel mehr verzweifelt genug, um endlich
aufzuwachen, meine Spiele zu durchschauen und aufzuhören, mich zu
belügen.»

Joh faltete die Hände zusammen und senkte die Augen.
Dann schaute er hoch und zeigte wieder sein Joh-Strahlen.

«Ich war Knall auf Fall gefangen von Manuela, weggepustet, komplett
verliebt. Sie war dann mindestens schwer zu halten, am Anfang, lange
glaubte ich immer wieder, sie zu verlieren. Mit ihr hatte sich etwas
verändert. Wie gesagt, davor war ich echt schlecht darin, eine gute
Beziehung zu halten, ein gebranntes Kind, nach ersten Hindernissen
schloss ich meine Gefühle weg. Ich war ein Versager und konnte mich
selbst kaum mehr aushalten.»

Louis runzelte die Stirn und beugte sich vor.

«Klingt hart.»

Joh räumte die Werkzeuge in die Kiste zurück und stand auf. Wieder umschloss er 
seinen verletzten Daumen.

«War es auch, allerspätestens nach meiner gescheiterten Ehe griff keine
meiner Lebensstrategien mehr. Meine Kinder hatte ich feige verlassen, unverzeihlich noch heute.»

Er wischte sich über die Augen und richtete seinen Blick direkt auf
Louis, bevor sich ein Anflug von Trauer zurückzog und sich sein Gesicht
aufhellte.

«Der Anfang mit Manuela war sehr eindrücklich,» er hielt inne, «weisst
du, ich war mir einfach nicht sicher, ob das was wird. Eines Morgens
schlurfte ich runter zu den Briefkästen unseres Hauses. Ich hatte mies
geschlafen, Mist geträumt und öffnete den Kasten.»

Joh streckte seine Arme von sich.

«Da sah ich das flache Päckchen mit Manuelas Absender. Ich war aufgeregt
wie ein Geburtstagskind, öffnete es noch unten im Stehen und starrte auf
eine CD in meinen Händen. Sie war von einem «Ivano Fossati», der
mir nichts sagte - was sich rasant ändern sollte,» er lachte herzhaft,
«ein Stück war speziell gekennzeichnet, \emph{«C’è tempo.»}

Joh unterbrach sich. Etwas Geheimnisvolles lag in der Luft.

«Interessiert es dich wirklich?»

«Na, klar.»

«Also, da lag eine weisse Karte bei, handschriftlich geschrieben.,» Joh
sprach schneller, als könne er es kaum erwarten, zum Finale zu kommen,
«darauf fand ich eine deutsche Übersetzung. Es sah aus, als müsste ich
verstehen, worum es ging. Kein Gruss, nichts. Nur aufgrund des Absenders
wusste ich, dass es Manuela war, die es mir geschickt hatte.»

Louis runzelte die Stirn.

«Zeit wofür? Klingt interessant. Den Song würde ich mir gern anhören.»

«Ja, machen wir, komm, wir haben eh eine Pause verdient, und ich freue mich
auch, den wieder mal zu hören.»

Joh winkte Louis hinter sich her vor die Hütte hinunter.

«Bin gleich wieder da.»

Louis setzte sich auf die Bank. Joh kam nach einer Weile zurück und
stellte einen Boomer auf den Tisch. Dann zog er sein Handy aus der
Hosentasche und suchte nach dem Lied.

«Das ist eine neue Aufnahme - aha, eine Live-Version, bin selbst
gespannt.»

Er wirkte erhitzt. Die Freude stand ihm ins Gesicht geschrieben. Sie
schauten sich an. Einvernehmlich, als wüssten sie beide genau, worum es ging. Louis’ Neugierde war vollends geweckt.
Und er konnte zusehen, wie Joh in sich verschwand.

\qrblock{Ivano Fossati \emph{«C’è tempo»}}{media/QR-Codes/042_C`è_tempo_(with_Ivano_Fossati)_Ivano_Fossati_Youtube.png}

«Phu, wow.»

Louis blinzelte ergriffen und nickte anerkennend.
Joh hatte ihn weit vorgelassen. Sein Gesicht sprach Bände. 

«Starker Text, und dann, wie, wie hast du reagiert?»

Joh holte tief Luft.

«Erstmal sass ich erschlagen auf meinem Sofa und wunderte mich. Ich
schwankte zwischen einem angenehmen Hochgefühl und grosser Verwirrung.
Ich fragte mich, was sie mir damit sagen wollte. War das ein Wink auf
gemeinsame Hoffnung? Ich fühlte mich im ersten Moment seltsam bedroht.
Vielleicht interpretierte ich den Text des Songs falsch? Meinte sie 
damit tatsächlich uns, mich?»

Louis fixierte ihn unentwegt und wartete ungeduldig auf die
Fortsetzung.

«Ja, meinte sie, jedenfalls machte mir die Geste ziemlich Angst. Ich
bewunderte sie grenzenlos. Darin lag der Haken. Schliesslich meldete ich
mich mehrere Tage nicht.»

Jetzt fing Joh zu Louis’ Verblüffung verschmitzt an zu lachen.

«Warum lachst du?»

Joh nahm sich die Brille ab und rieb die Gläser mit einem Taschentuch sauber. Er war wieder am Denken. Es konnte
interessant werden.

«Stell dir vor, Ciro, du lernst eine Frau kennen, die dir mehr als
gefällt, hast aber einfach Schiss, dich wirklich, hm, ganz einzulassen,»
er rückte sich in eine angenehmere Position, «ich ahnte damals schon,
dass Manuela keine Frau für eine lose Liebschaft war, und dann stell dir
vor, in der Vergangenheit hast du dich einfach davongemacht, wenn du es
am wenigsten wolltest, und dann schickt sie dir sowas. Was hättest
\emph{du} getan?»

Louis zuckte perplex zusammen. Joh lehnte sich nach hinten zurück und
setzte sich die Brille wieder auf. Louis glaubte, wieder seinen
eindringlichen Blick zu lesen. 
Die Frage sprengte sein Wohlgefühl. Louis spitzte die Lippen, hob die
Schultern und fühlte eine unangenehme Hitze hochkommen 
Joh wischte sich kurz und nachdenklich mit der Hand über den Mund, stand
auf und klopfte Louis auf die Schulter.

«Dämlich von mir,» er lächelte, «verzeih, die Antworten musste ich
natürlich bei mir finden,» dann drehte er sich nachdenklich zu Louis,
«was ich damals zu verstehen lernte, \emph{Alles hat seine Zeit},
Manuela hatte viel und gab mir davon ab. Mit «Fossatis» Song
wollte sie mir den Schlüssel reichen. - Also dann, ich mach mich ans
Kochen. Ich würde mich freuen, wenn du mir hilfst.»

Johs Handy klingelte. Er klaubte es aus seiner Jackentasche.

«Klar,» sagte Louis schnell und seine Anspannung löste sich.

Bevor Joh den Anruf entgegennahm, drehte er sich beim Weggehen noch einmal zu Louis
um.

«Wenn du willst, erzähle ich dir den Rest unserer Anfangsgeschichte beim
Kochen?»

«Ja, gerne.»

Johs Frage freute Louis. 
Er hatte sich in etwas hineingesteigert. Liebesbeziehungen hatten es an
sich, kompliziert zu sein. Er war nicht der einzige, der sich schwer
tat. Eigentlich tröstete ihn der Gedanke, dass auch dieser scheinbar so abgeklärte Kerl seine Kämpfe
ausgefochten hatte. 
Joh beendete den Anruf. Er legte das Handy auf den Esstisch.

«Schade, Ciro, David kommt unerwartet früher zurück, meine Geschichte
muss warten,» er wackelte mit dem Kopf, «wir finden bestimmt später
einen ruhigen Moment zusammen.»

Louis nickte. Er liess sich seine Enttäuschung nicht anmerken.
Schon stand David in der Eingangstür. Sie kochten gemeinsam. Dabei
redeten sie vorwiegend und ausführlich über verschiedene
Zubereitungsarten von Gemüseaufläufen.
